% GENERATED FILE. Edit canonical lessons, then run npm run generate:books.
% canonical-source-hash: fnv1a64:f98875afde60b8c3
% canonical-lessons: ZH-C107-tiao, ZH-C107-pa-climb, ZH-C107-ban, ZH-C107-ting-stop, ZH-C107-likai

\chapter{On the Move}
\label{ch:zh-on-the-move}
\hlchaptermodality{107}

\begin{chapteropening}
\textbf{By the end of this chapter:} \emph{I can say jump, climb, carry, stop and leave.}

Complete the last lesson of chapter 107: I can say jump, climb, carry, stop and leave.
\end{chapteropening}

\section[tiào]{tiào — to jump}

\label{lesson:ZH-C107-tiao}



\begin{quote}
Before the new one: say the Mandarin for bicycle, then the Mandarin for to ride.
\end{quote}

\subsection*{You'll want to know: tiào}
\emph{tiào} — \textquotedblleft{}to jump\textquotedblright{}.

\begin{quote}
\emph{wǒ de gǒu hěn xǐhuan tiào} — My dog loves jumping.
\end{quote}

Its characters come later; for now it is a word for the ear.

\begin{cousinweb}
\emph{tiào} is inherited from Middle Chinese \emph{dew} (Guangyun, in Baxter's transcription), glossed 'to leap'. The regular change from \emph{dew} gives tiáo, with a rising tone, not the falling tiào of today; how it came to its tone is not settled.
\end{cousinweb}

\begin{grammarlens}[title={What you've built}]
One more word for everyday actions.
\end{grammarlens}

\subsection*{Your turn}
\begin{itemize}
\raggedright
  \item \emph{Say it:} \emph{tiào}
  \item \emph{Say it:} \emph{tiào}, once more
  \item \emph{Say it:} \emph{wǒ de gǒu hěn xǐhuan tiào} — My dog loves jumping.
  \item \emph{Recall:} say the Mandarin for bicycle, then the Mandarin for to ride, then say \emph{tiào} again
\end{itemize}

\begin{tcolorbox}[breakable,colback=teal!4,colframe=teal!35!black,title={Before you move on}]
What does tiào mean? (\textquotedblleft{}To jump\textquotedblright{}.) Say it once more.
\end{tcolorbox}
\section[pá]{pá — to climb; to crawl}

\label{lesson:ZH-C107-pa-climb}



\begin{quote}
Before the new one: say the Mandarin for to ride, then the Mandarin for to jump.
\end{quote}

\subsection*{You'll want to know: pá}
\emph{pá} — \textquotedblleft{}to climb; to crawl\textquotedblright{}.

\begin{quote}
\emph{wǒmen míngtiān qù pá shān} — We're going up into the hills tomorrow.
\end{quote}

Its characters come later; for now it is a word for the ear.

\begin{cousinweb}
\emph{pá} is inherited from Middle Chinese \emph{bae} (Guangyun, in Baxter's transcription), glossed 'to scratch'; how it came to mean to climb is not settled.
\end{cousinweb}

\begin{grammarlens}[title={What you've built}]
One more word for everyday actions.
\end{grammarlens}

\subsection*{Your turn}
\begin{itemize}
\raggedright
  \item \emph{Say it:} \emph{pá}
  \item \emph{Say it:} \emph{pá}, once more
  \item \emph{Say it:} \emph{wǒmen míngtiān qù pá shān} — We're going up into the hills tomorrow.
  \item \emph{Recall:} say the Mandarin for to ride, then the Mandarin for to jump, then say \emph{pá} again
\end{itemize}

\begin{tcolorbox}[breakable,colback=teal!4,colframe=teal!35!black,title={Before you move on}]
What does pá mean? (\textquotedblleft{}To climb; to crawl\textquotedblright{}.) Say it once more.
\end{tcolorbox}
\section[bān]{bān — to move house; to move (something heavy)}

\label{lesson:ZH-C107-ban}



\begin{quote}
Before the new one: say the Mandarin for to jump, then the Mandarin for to climb.
\end{quote}

\subsection*{You'll want to know: bān}
\emph{bān} — \textquotedblleft{}to move house; to move (something heavy)\textquotedblright{}.

\begin{quote}
\emph{nǐ néng bāng wǒ bān zhuōzi ma} — Can you help me move the table?
\end{quote}

Its characters come later; for now it is a word for the ear.

\begin{cousinweb}
\emph{bān} has no older form on record under this character. The dictionary lists the same sound, \emph{pan} (Guangyun, in Baxter's transcription), under an older character glossed 'to transport', which is probably the same word, though that is not settled.
\end{cousinweb}

\begin{grammarlens}[title={What you've built}]
One more word for everyday actions.
\end{grammarlens}

\subsection*{Your turn}
\begin{itemize}
\raggedright
  \item \emph{Say it:} \emph{bān}
  \item \emph{Say it:} \emph{bān}, once more
  \item \emph{Say it:} \emph{nǐ néng bāng wǒ bān zhuōzi ma} — Can you help me move the table?
  \item \emph{Recall:} say the Mandarin for to jump, then the Mandarin for to climb, then say \emph{bān} again
\end{itemize}

\begin{tcolorbox}[breakable,colback=teal!4,colframe=teal!35!black,title={Before you move on}]
What does bān mean? (\textquotedblleft{}To move house; to move (something heavy)\textquotedblright{}.) Say it once more.
\end{tcolorbox}
\section[tíng]{tíng — to stop; to park}

\label{lesson:ZH-C107-ting-stop}



\begin{quote}
Before the new one: say the Mandarin for to climb, then the Mandarin for to move house.
\end{quote}

\subsection*{You'll want to know: tíng}
\emph{tíng} — \textquotedblleft{}to stop; to park\textquotedblright{}.

\begin{quote}
\emph{chē tíng le} — The car has stopped.
\end{quote}

Its characters come later; for now it is a word for the ear.

\begin{cousinweb}
\emph{tíng} is inherited from Middle Chinese \emph{deng} (Guangyun, in Baxter's transcription), glossed 'to rest', 'to settle' and 'to stop'.
\end{cousinweb}

\begin{grammarlens}[title={What you've built}]
One more word for everyday actions.
\end{grammarlens}

\subsection*{Your turn}
\begin{itemize}
\raggedright
  \item \emph{Say it:} \emph{tíng}
  \item \emph{Say it:} \emph{tíng}, once more
  \item \emph{Say it:} \emph{chē tíng le} — The car has stopped.
  \item \emph{Recall:} say the Mandarin for to climb, then the Mandarin for to move house, then say \emph{tíng} again
\end{itemize}

\begin{tcolorbox}[breakable,colback=teal!4,colframe=teal!35!black,title={Before you move on}]
What does tíng mean? (\textquotedblleft{}To stop; to park\textquotedblright{}.) Say it once more.
\end{tcolorbox}
\section[líkāi]{líkāi — to leave a place}

\label{lesson:ZH-C107-likai}



\begin{quote}
Before the new one: say the Mandarin for to move house, then the Mandarin for to stop.
\end{quote}

\subsection*{You'll want to know: líkāi}
\emph{líkāi} — \textquotedblleft{}to leave a place\textquotedblright{}.

\begin{quote}
\emph{wǒmen jǐ diǎn líkāi} — What time are we leaving?
\end{quote}

Its characters come later; for now it is a word for the ear.

\begin{cousinweb}
\emph{líkāi} is built from lí 'to part' and \zh{开} 'to open', both inherited from Middle Chinese, \emph{lje} and \emph{khoj} (Guangyun, in Baxter's transcription); the dictionary says parting at a short distance is \emph{lje}, at a long one is another word.
\end{cousinweb}

\begin{grammarlens}[title={What you've built}]
One more word for everyday actions.
\end{grammarlens}

\subsection*{Your turn}
\begin{itemize}
\raggedright
  \item \emph{Say it:} \emph{líkāi}
  \item \emph{Say it:} \emph{líkāi}, once more
  \item \emph{Say it:} \emph{wǒmen jǐ diǎn líkāi} — What time are we leaving?
  \item \emph{Recall:} say the Mandarin for to move house, then the Mandarin for to stop, then say \emph{líkāi} again
\end{itemize}

\begin{tcolorbox}[breakable,colback=teal!4,colframe=teal!35!black,title={Before you move on}]
What does líkāi mean? (\textquotedblleft{}To leave a place\textquotedblright{}.) Say it once more.
\end{tcolorbox}
